\begin{figure}[ht]
\centering
\begin{tikzpicture}[
  pcircle/.style = {draw=engblue, very thick},
  bit/.style     = {circle, draw=engblack, thick, fill=englime!25,
                    inner sep=0pt, minimum size=6.5mm, font=\small},
  par/.style     = {circle, draw=engred, thick, fill=engred!12,
                    inner sep=0pt, minimum size=6.5mm, font=\small},
]
% The classic three-circle Venn diagram for Hamming(7,4).
% Parity bits p1,p2,p3 each cover the data bits in its circle.
\draw[pcircle] (-1.05, 0.55) circle (2cm);
\draw[pcircle] (1.05, 0.55)  circle (2cm);
\draw[pcircle] (0, -1.2)     circle (2cm);

% data bits d1..d4 in the overlap regions
\node[bit] (d1) at (0, 1.15)      {$d_1$};
\node[bit] (d2) at (-0.95, -0.35) {$d_2$};
\node[bit] (d3) at (0.95, -0.35)  {$d_3$};
\node[bit] (d4) at (0, -0.55)     {$d_4$};

% parity bits p1,p2,p3 in the single-circle regions
\node[par] (p1) at (-1.7, 1.45)  {$p_1$};
\node[par] (p2) at (1.7, 1.45)   {$p_2$};
\node[par] (p3) at (0, -2.35)    {$p_3$};

\node[engblue, font=\footnotesize, anchor=east] at (-2.7, 1.9) {circle $1$};
\node[engblue, font=\footnotesize, anchor=west] at (2.7, 1.9)  {circle $2$};
\node[engblue, font=\footnotesize, anchor=north] at (0, -3.25) {circle $3$};
\end{tikzpicture}
\caption{The Hamming$(7,4)$ code as three overlapping parity
circles. Four data bits $d_1, d_2, d_3, d_4$ occupy the overlap
regions; each parity bit $p_j$ is set so that its circle holds an
even number of ones. A single bit-flip makes exactly the circles
containing it fail their parity check, and the failing-circle
pattern is the binary address of the corrupted bit, so one error is
located and corrected. The placement $d_4$ in all three circles and
each $p_j$ in one circle realises the parity-check structure of the
worked example.}
\label{fig:vol02:ch17:hamming}
\end{figure}
